% se puede agregar la opción [english] para 
%  memorias o tesis en inglés (borrando el archivo .aux)
\documentclass{tesisUOH} 
\usepackage{subfiles}
\usepackage{fancyhdr}

\author{Nombre Completo Autor}
\title{\textbf{Título de la Tesis/Memoria}}

% incluir ambos comandos para una doble titulación
%  o quitar el comando que no aplica
\memoria{Ingenier[o/a?] Civil en ???}
\tesis{Magíster en Ciencias de la Ingeniería,\\mención ????}
%\tesis{Doctor en ???} % incluir solo este comando para doctorados

% puede haber varios profesores guía seperados por coma;
% pero si es una memoria, solo puede haber un profesor guía
\guia{Nombre Completo Guía} 

% puede haber varios profesores co-guía seperados por coma;
% pero si es una memoria, el profesor co-guía será el primer
% integrante de la comisión
%\coguia{Nombre Completo Co-Guía} % incluir en caso de co-guía de *tesis*

%\cotutela{Nombre Institución} % incluir en caso de cotutela

\comision{Nombre Completo Uno,Nombre Completo Dos,Nombre Completo Tres}

%\auspicio{Nombre Institución} % incluir en caso de recibir financiamiento

% tiene que ser el año en que se da el examen de título/grado (defensa)
%\anho{2021} % incluir solo para reemplazar el año actual

\usepackage{lipsum}

\begin{document}

\frontmatter
\maketitle

\begin{resumen}
\lipsum[1-4]
\end{resumen}

% opcional: incluir para tesis en inglés;
%  en este caso hay que tener el resumen y abstract
%   en ambos idiomas
%\begin{abstract}
%\lipsum[1-4]
%\end{abstract}

\begin{dedicatoria}
Una dedicatoria corta.
\end{dedicatoria}

\begin{thanks}
\lipsum[1-2]
\end{thanks}
{
\hypersetup{linkcolor=black}
\tableofcontents
\listoftables % opcional
\listoffigures % opcional
}

\pagestyle{fancy}
\fancyhead{} % clear all header fields
\fancyhead[C]{\nouppercase{\rightmark\hfill\leftmark}}
\fancyhead[C]{\nouppercase{\leftmark\hfill\rightmark}}
\fancyfoot{} % clear all footer fields
\fancyfoot[C]{\thepage}

\mainmatter

\subfile{cap1}
\subfile{cap2}
\subfile{cap3}
\subfile{conclu}

%Inlcuir cuantos capítulos estime necesario
% \input{cap1.tex}
% \input{cap2.tex}
% \input{cap3.tex}
% \input{conclu.tex}

% ver https://www.overleaf.com/learn/latex/Glossaries
% \input{glosario.tex} % opcional

%\nocite{*} %para incluir todos los items de la bibliografia (aun cuando no sean citados) 
\bibliographystyle{plain}
\bibliography{bibliografia}

% opcional ...
\begin{appendices}
\subfile{anexoA}
\end{appendices}
\end{document}